\documentclass[11pt]{article}
\usepackage[margin=1in]{geometry}
\usepackage{booktabs,graphicx,amsmath,hyperref,xcolor}
\usepackage[numbers]{natbib}
\IfFileExists{svg.sty}{\usepackage{svg}}{}

% Results: run `bb report runs/experiments/<run>` on the GPU machine, then copy that run's
% figures/ and tables/ directories next to this file. \result falls back to a visible note.
\newcommand{\result}[2]{\IfFileExists{#1}{#2}{\fbox{\parbox{0.9\linewidth}{\small Missing #1: run \texttt{bb report} and copy its output here.}}}}

\title{Intelligible Blackboard Architecture with Counterfactual Agents}
\author{Team (see docs/team.md)}
\date{}

\begin{document}
\maketitle

\begin{abstract}
Agents that pass text down a chain lose context and compound each other's errors. We let three
language-model agents coordinate only through a shared blackboard on which every message follows
the Two-Way Intelligibility Protocol (PXP): a prediction, an explanation, and a tag derived from
the sender's judgments of the message it answers. When the board deadlocks, agents with
counterfactual capability replay the session in an isolated sandbox with alternatives to their
own past messages, assign credit by how much closer each alternative brings the team to
consensus, and apply the best one to the live board through a lesson in their prompt, a revised
position and a lower voting weight. We vary the share of counterfactual-capable agents (0, 33,
67, 100\%) over 500 tasks from MSCoRe, KramaBench and MedAgentBench and report convergence
success, intelligibility depth and token cost.
\end{abstract}

\section{Introduction}
% Motivation: chains vs a shared structured state; PXP; deadlocks; counterfactual credit.

\section{PXP on a blackboard}
\label{sec:pxp}
Every message carries the sender's prediction $y$ and explanation $e$. Following Algorithm~1 of
\citet{mestha2025}, the sender compares the peer message $(y', e')$ (the latest message from
another agent) with its own previous position $(y'', e'')$ and its new output with that position:
\[
\textsc{match} = \mathrm{MATCH}(y', y''),\quad \textsc{agree} = \mathrm{AGREE}(e', e''),\quad
\textsc{changed} = \neg\mathrm{MATCH}(y, y'') \lor \neg\mathrm{AGREE}(e, e'').
\]
Message 1 is \textsc{init}; \textsc{match}$\land$\textsc{agree} gives \textsc{ratify}; exactly
one of them gives \textsc{revise} if \textsc{changed} and \textsc{refute} otherwise; neither gives
\textsc{reject} after message $k$ and \textsc{revise}/\textsc{refute} before. The model supplies
the judgments that need understanding; the agent's procedure computes MATCH over listed options
and the tag; a grammar checks every message on commit and rejects (never repairs) invalid ones.
Session intelligibility follows Definitions~1--2 of \citet{mestha2025}, taking the weakest level
over agents (Section~\ref{sec:metrics}).

\section{Architecture}
% Four layers (docs/architecture.md): event-sourced blackboard with optimistic versioning;
% event-driven scheduler with a priority arbiter and weighted consensus over RATIFY chains;
% agents and providers (vLLM, constrained decoding); counterfactual layer; UI.

\section{Counterfactual credit assignment}
On a deadlock, each capable agent reflects on its last three messages, proposes up to three
grammatical alternatives per message, and replays the session from each in a sandbox rebuilt from
the event log. The credit of an alternative is the gain in consensus proximity (the weighted
share of the largest agreeing group) between the live board at the deadlock and the end of the
replay. The best alternative with positive credit is applied: its lesson joins the agent's system
prompt, the agent's position is revised, and its weight becomes $\max(0.25,\, w(1-0.5\cdot
\mathrm{credit}))$; the session resumes with the agent posting next.

\section{Experimental setup}
\label{sec:metrics}
% Copy from docs/experiments.md: groups, benchmarks and their adaptation, model, metrics.

\section{Results}
\begin{table}[h]
\centering\small
\result{tables/groups.tex}{\input{tables/groups.tex}}
\caption{Results per group and benchmark. Convergence success with 95\% Wilson intervals.}
\end{table}

\begin{figure}[h]
\centering
\result{figures/convergence.png}{\includegraphics[width=0.75\linewidth]{figures/convergence.png}}
\caption{Convergence success against the share of counterfactual-capable agents.}
\end{figure}

\begin{figure}[h]
\centering
\result{figures/intelligibility.png}{\includegraphics[width=0.75\linewidth]{figures/intelligibility.png}}
\caption{Intelligibility depth of whole sessions.}
\end{figure}

\begin{figure}[h]
\centering
\result{figures/tokens.png}{\includegraphics[width=0.75\linewidth]{figures/tokens.png}}
\caption{Tokens per instance on the live board and in the counterfactual sandbox.}
\end{figure}

\section{Discussion and limitations}
% See docs/experiments.md#limitations.

\bibliographystyle{plainnat}
\bibliography{references}
\end{document}
